% =========================================================================
% TEOREMA DO FLUXO MÁXIMO E CORTE MÍNIMO (MFMC)
% =========================================================================
\chapter{O Teorema Fundamental do Fluxo Máximo e Corte Mínimo}\label{sec:mfmc_capitulo}

\section{Redes Residuais, Caminhos Aumentantes e Dualidade}

Para determinar se um fluxo pode ser aumentado e permitir a reversão de decisões de roteamento anteriores, introduz-se o conceito de rede residual (veja~\cite{ahuja1993network, ford1956maximal}).

\begin{definicao}[Capacidade Residual e Digrafo Residual][def:residual]
Dado um fluxo viável $f$ operando sobre uma rede $D = (V, A)$ com capacidades originais $c: A \to \mathbb{R}_{\ge 0}$:
\begin{enumerate}
    \item A \textbf{capacidade residual} de qualquer par ordenado $(u,v) \in V \times V$, denotada por $c_f(u,v)$, é definida por:
    \begin{equation}
        c_f(u,v) \coloneqq c(u,v) - f(u,v) + f(v,u)
    \end{equation}
    convencionando-se $c(u,v) = 0$ e $f(u,v) = 0$ caso $(u,v) \notin A$.
    \item O \textbf{digrafo residual} $D_f = (V, A_f)$ é o grafo orientado cujos arcos são aqueles com capacidade residual estritamente positiva:
    \begin{equation}
        A_f \coloneqq \{ (u,v) \in V \times V \mid c_f(u,v) > 0 \}
    \end{equation}
\end{enumerate}
O termo $+ f(v,u)$ representa a capacidade residual reversa, permitindo cancelar fluxo enviado no sentido oposto.
\end{definicao}

Um caminho aumentante em $D_f$ é um caminho simples direcionado com origem na fonte $s$ e término no sumidouro $t$. A existência de um caminho aumentante $P$ implica que o fluxo $|f|$ pode ser incrementado por $\delta = \min_{(u,v) \in P} c_f(u,v) > 0$, denominado capacidade residual gargalo de $P$.

A Figura~\ref{fig:rede_residual} ilustra uma rede com fluxo parcial e seu correspondente digrafo residual $D_f$, evidenciando os arcos diretos (capacidade ociosa) e inversos (fluxo cancelável).

\begin{figure}[!ht]
\centering
\begin{tikzpicture}[thick]

\begin{scope}[xshift=0cm]
  \node[source node] (s) at (0,1.2) {$s$};
  \node[flow node] (a) at (2,2.4) {$a$};
  \node[flow node] (b) at (2,0) {$b$};
  \node[sink node] (t) at (4,1.2) {$t$};
  \draw[flow edge] (s) -- node[edge lbl above] {$2/3$} (a);
  \draw[flow edge] (s) -- node[edge lbl below] {$1/2$} (b);
  \draw[flow edge] (a) -- node[edge lbl, right] {$1/1$} (b);
  \draw[flow edge] (a) -- node[edge lbl above] {$1/3$} (t);
  \draw[flow edge] (b) -- node[edge lbl below] {$2/3$} (t);
  \node[font=\footnotesize,below] at (2,-0.7) {\textit{(a) Rede original} $f/c$};
\end{scope}

\node at (5.2,1.2) {\Huge $\Rightarrow$};

\begin{scope}[xshift=6.5cm]
  \node[source node] (s2) at (0,1.2) {$s$};
  \node[flow node] (a2) at (2,2.4) {$a$};
  \node[flow node] (b2) at (2,0) {$b$};
  \node[sink node] (t2) at (4,1.2) {$t$};

  \draw[flow edge,blue] (s2) -- node[edge lbl below, text=blue!80!black] {$1$} (a2);
  \draw[flow edge,blue] (s2) -- node[edge lbl above, text=blue!80!black] {$1$} (b2);
  \draw[flow edge,blue] (a2) -- node[edge lbl below, text=blue!80!black] {$2$} (t2);
  \draw[flow edge,blue] (b2) -- node[edge lbl above, text=blue!80!black] {$1$} (t2);

  \draw[rev edge,bend right=30] (a2) to node[edge lbl above, text=red!80!black] {$2$} (s2);
  \draw[rev edge,bend left=30] (b2) to node[edge lbl below, text=red!80!black] {$1$} (s2);
  \draw[rev edge] (b2) -- node[edge lbl, left, text=red!80!black] {$1$} (a2);
  \draw[rev edge,bend right=30] (t2) to node[edge lbl above, text=red!80!black] {$1$} (a2);
  \draw[rev edge,bend left=30] (t2) to node[edge lbl below, text=red!80!black] {$2$} (b2);
  \node[font=\footnotesize,below] at (2,-0.7) {\textit{(b) Digrafo residual} $D_f$};
\end{scope}
\end{tikzpicture}
\caption{Rede com fluxo parcial (esquerda, valores $f/c$) e o digrafo residual $D_f$ resultante (direita). Arcos azuis sólidos: capacidade ociosa ($c_f > 0$). Arcos vermelhos tracejados: capacidade reversa (fluxo cancelável).}
\label{fig:rede_residual}
\end{figure}
\FloatBarrier

O critério de parada e o certificado de otimalidade para fluxos apoiam-se no conceito de corte:

\begin{definicao}[Corte $s-t$ e Capacidade de Corte][def:corte]
Dada uma rede de fluxo $D = (V, A)$ com capacidades $c: A \to \mathbb{R}_{\ge 0}$, fonte $s$ e sumidouro $t$:
\begin{enumerate}
    \item Um \textbf{corte $s-t$} é uma partição do conjunto de vértices $V$ em dois subconjuntos disjuntos $S$ e $T = V \setminus S$ tais que $s \in S$ e $t \in T$.
    \item A \textbf{capacidade do corte}, denotada por $\ccap(S, T)$, é a soma das capacidades dos arcos direcionados com cauda em $S$ e cabeça em $T$:
    \begin{equation}
        \ccap(S, T) \coloneqq \sum_{(u,v) \in A, \, u \in S, \, v \in T} c(u,v)
    \end{equation}
\end{enumerate}
\textbf{Observação:} Arcos de $T$ para $S$ não entram no cálculo de $\ccap(S, T)$, pois cruzam no sentido oposto ao corte.
\end{definicao}

O Teorema do Fluxo Máximo e Corte Mínimo estabelece a equivalência entre o valor ótimo do fluxo e a capacidade mínima de corte:

\begin{teorema}[Fluxo Máximo e Corte Mínimo --- Ford e Fulkerson, 1956][teo:mfmc]
Em qualquer rede de fluxo $D = (V, A)$ com capacidades não negativas, fonte $s$ e sumidouro $t$, o valor do fluxo máximo é igual à capacidade do corte $s-t$ mínimo:
\begin{equation}
    \max_f |f| = \min_{(S,T)} \ccap(S,T)
\end{equation}
Além disso, um fluxo viável $f$ é máximo se, e somente se, o digrafo residual $D_f$ não contém nenhum caminho aumentante de $s$ a $t$.
\end{teorema}

\intuicao{Princípio do Gargalo} A relação entre fluxo e corte fundamenta-se no princípio do gargalo: qualquer partição $(S, T)$ que separe a fonte $s$ do sumidouro $t$ impõe um limite superior sobre a quantidade de fluxo transportável, pois todo o fluxo líquido que transita de $s$ a $t$ deve cruzar os arcos direcionados de $S$ para $T$. Deste modo, $|f| \le \ccap(S, T)$ para qualquer fluxo viável $f$ e qualquer corte $s-t$. O teorema afirma que o valor máximo do fluxo coincide com a capacidade mínima de um corte: quando o fluxo atinge seu valor máximo, os arcos que cruzam o corte mínimo encontram-se saturados ($f(u,v) = c(u,v)$), inviabilizando caminhos aumentantes na rede residual.


\begin{demonstracao}
A demonstração deste teorema é construtiva e baseia-se na estrutura do digrafo residual (veja~\cite{ford1956maximal, cook1998combinatorial}). Pelo Lema da Capacidade do Corte, $|f| \le \ccap(S, T)$ para qualquer fluxo viável $f$ e qualquer partição $s-t$. Se existir um fluxo $f$ e um corte $(S, T)$ tais que $|f| = \ccap(S, T)$, então $f$ é um fluxo máximo e $(S, T)$ é um corte mínimo.

Suponha que $f$ seja um fluxo máximo. Se existisse um caminho aumentante $P$ de $s$ a $t$ em $D_f$, seria possível definir um fluxo $f'$ com $|f'| = |f| + c_f(P) > |f|$, contradizendo a maximalidade de $f$. Logo, não existe caminho aumentante em $D_f$ (veja~\cite{ahuja1993network}). Definimos o conjunto $S$ como o conjunto de todos os vértices que são alcançáveis a partir da fonte $s$ em $D_f$ através de arcos com capacidade residual estritamente positiva ($c_f > 0$). Definimos $T = V \setminus S$ como o restante dos vértices.

Como não há caminho aumentante, o sumidouro $t$ é inalcançável a partir de $s$ em $D_f$, logo $t \in T$ e $s \in S$. Temos, portanto, um corte $s-t$ válido particionando a rede.

\begin{figure}[!ht]
\centering
\begin{tikzpicture}[thick]
    \draw[cut S] (-0.8, -1.8) rectangle (2.8, 1.8);
    \node[cut label S] at (-0.6, 1.65) {$\mathbf{S}$};

    \draw[cut T] (5.4, -1.8) rectangle (9.0, 1.8);
    \node[cut label T] at (8.8, 1.65) {$\mathbf{T}$};

    \node[source node] (s) at (0.0, 0.0) {$s$};
    \node[s node internal] (a) at (2.0, 0.8) {$a$};
    \node[s node internal] (b) at (2.0, -0.8) {$b$};

    \node[t node internal] (c) at (6.2, 0.8) {$c$};
    \node[t node internal] (d) at (6.2, -0.8) {$d$};
    \node[sink node] (t) at (8.2, 0.0) {$t$};

    \draw[line width=1.5pt, dashed, cutviolet] (4.1, -1.85) -- (4.1, 1.85)
        node[above=2.5pt, solid, font=\scriptsize\bfseries, text=cutviolet!90!black, fill=white, inner sep=2.5pt, rounded corners=3pt, draw=cutviolet!70] {Corte Mínimo $(S,T)$};

    \draw[sat edge] (a) -- node[edge lbl above, text=cutviolet!90!black] {$f=c$} (c);
    \draw[sat edge] (b) -- node[edge lbl below, text=cutviolet!90!black] {$f=c$} (d);

    \draw[flow edge] (s) -- (a);
    \draw[flow edge] (s) -- (b);
    \draw[flow edge] (c) -- (t);
    \draw[flow edge] (d) -- (t);
    \draw[flow edge] (a) -- (b);
    \draw[flow edge] (c) -- (d);
\end{tikzpicture}

\caption{Corte mínimo $(S,T)$: os arcos que cruzam de $S$ para $T$ (em violeta) estão saturados ($f = c$).}
\label{fig:corte_minimo}
\end{figure}

Para qualquer arco original $(u,v) \in A$ tal que $u \in S$ e $v \in T$, sabemos que $v$ não é alcançável a partir de $u$ em $D_f$. Isto implica obrigatoriamente que a capacidade residual $c_f(u,v) = 0$. Pela definição de capacidade residual, isto só ocorre se o arco original estiver completamente saturado, ou seja, $f(u,v) = c(u,v)$. De forma análoga, para qualquer arco no sentido reverso, de $T$ para $S$, o fluxo deve ser nulo ($f(v,u) = 0$), caso contrário haveria capacidade residual reversa permitindo a travessia de $S$ para $T$ (veja~\cite{cook1998combinatorial}).

Conclui-se que o fluxo líquido através do corte $(S, T)$ é exatamente a soma das capacidades dos arcos de $S$ para $T$, o que corresponde à capacidade total do corte. Como o fluxo da rede deve passar através deste gargalo, provamos que o valor $|f|$ iguala a capacidade deste corte, atestando a otimalidade de ambos (veja~\cite{ford1956maximal}). \end{demonstracao}

\implicacoes{e Certificado de Otimalidade} O Teorema Max-Flow Min-Cut fornece um certificado de dualidade forte de otimalidade. Enquanto a busca exaustiva por cortes mínimos exigiria avaliar $2^{|V|-2}$ partições bipartidas, o teorema assegura que o corte mínimo pode ser obtido em tempo linear $\mathcal{O}(|V| + |A|)$ executando uma busca em largura (BFS) a partir de $s$ sobre a rede residual final produzida por qualquer algoritmo de fluxo máximo (como Edmonds-Karp ou Dinic). Além disso, este teorema permite demonstrar teoremas estruturais em teoria dos grafos, incluindo König, Hall e Menger.


\section{A Formulação Primal-Dual em Programação Linear}\label{sec:dualidade_primal_dual}

A equivalência entre o valor do fluxo máximo e a capacidade do corte mínimo admite demonstração algébrica sob a ótica da Programação Linear~\cite{cook1998combinatorial, schrijver2003combinatorial}. O Teorema do Fluxo Máximo e Corte Mínimo corresponde à dualidade forte de von Neumann e Dantzig aplicada a redes direcionadas.

O Problema do Fluxo Máximo consiste em maximizar o valor do fluxo $|f|$, formulado como o seguinte Programa Linear Primal:
\begin{align}
    \max \quad & \sum_{(s,v) \in A} f(s,v) - \sum_{(u,s) \in A} f(u,s) \label{eq:primal_obj} \\[4pt]
    \sa & \sum_{(v,w) \in A} f(v,w) - \sum_{(u,v) \in A} f(u,v) = 0, \quad \forall v \in V \setminus \{s, t\} \label{eq:primal_conservacao} \\[4pt]
    & 0 \le f(u,v) \le c(u,v), \quad \forall (u,v) \in A \label{eq:primal_capacidade}
\end{align}

Para derivar o Programa Linear Dual correspondente, associam-se multiplicadores de Lagrange (variáveis duais) às restrições do primal:
\begin{enumerate}
    \item Para cada restrição de conservação de fluxo no vértice $v \in V \setminus \{s,t\}$ (Equação~\ref{eq:primal_conservacao}), introduz-se uma variável dual irrestrita $p(v) \in \mathbb{R}$, interpretada como um \textbf{potencial nodal}. Convenciona-se $p(s) = 1$ e $p(t) = 0$.
    \item Para cada restrição de capacidade superior $f(u,v) \le c(u,v)$ (Equação~\ref{eq:primal_capacidade}), introduz-se uma variável dual não negativa $d(u,v) \ge 0$, que quantifica a \textbf{tensão} ao longo do arco $(u,v)$.
\end{enumerate}

Pela teoria da dualidade linear, a maximização do fluxo converte-se na minimização do produto das capacidades pelas tensões dos arcos. O Programa Linear Dual resultante assume a forma:
\begin{align}
    \min \quad & \sum_{(u,v) \in A} c(u,v) \cdot d(u,v) \label{eq:dual_obj} \\[4pt]
    \sa & d(u,v) - p(u) + p(v) \ge 0, \quad \forall (u,v) \in A \label{eq:dual_tensao} \\[4pt]
    & d(u,v) \ge 0, \quad \forall (u,v) \in A \\[4pt]
    & p(s) = 1, \quad p(t) = 0, \quad p(v) \in \mathbb{R}, \quad \forall v \in V
\end{align}

Para qualquer arco $(u,v)$, tem-se $d(u,v) \ge p(u) - p(v)$. Como $c(u,v) \ge 0$, a solução ótima dual satisfaz:
\begin{equation}
    d(u,v) = \max\{0, \, p(u) - p(v)\}
\end{equation}

Em um programa linear arbitrário, as variáveis duais $p(v)$ poderiam assumir valores fracionários em $[0, 1]$. Contudo, a matriz de coeficientes do sistema de restrições é a transposta da matriz de incidência do digrafo $D$. Conforme demonstrado na Subseção~\ref{sec:tum_secao}, a matriz de incidência de um digrafo é \textbf{Totalmente Unimodular (TUM)}.

Pela teoria dos Poliedros e da Programação Inteira~\cite{cook1998combinatorial, korte2018combinatorial}, a total unimodularidade da matriz de restrições garante que todos os vértices do poliedro dual são inteiros. Consequentemente, para capacidades inteiras, existe solução básica ótima dual onde os potenciais nodais são binários:
\begin{equation}
    p^*(v) \in \{0, 1\}, \quad \forall v \in V
\end{equation}

Essa propriedade binária induz uma partição de corte:
\begin{equation}
    S = \{v \in V \mid p^*(v) = 1\} \quad \text{e} \quad T = \{v \in V \mid p^*(v) = 0\}
\end{equation}
Como $p^*(s) = 1$ e $p^*(t) = 0$, tem-se $s \in S$ e $t \in T$, definindo um corte $s-t$.

Sob esta partição binária, as variáveis $d^*(u,v) = \max\{0, \, p^*(u) - p^*(v)\}$ resultam em:
\begin{equation}
    d^*(u,v) = \begin{cases}
        1, & \text{se } u \in S \text{ e } v \in T; \\
        0, & \text{caso contrário}.
    \end{cases}
\end{equation}
Substituindo $d^*(u,v)$ na função objetivo dual~\eqref{eq:dual_obj}:
\begin{equation}
    \sum_{(u,v) \in A} c(u,v) \cdot d^*(u,v) = \sum_{u \in S, \, v \in T} c(u,v) = \ccap(S, T)
\end{equation}

Pelo Teorema da Dualidade Forte:
\begin{equation}
    \max |f| = \min \sum_{(u,v) \in A} c(u,v) d(u,v) = \min_{(S,T)} \ccap(S,T)
\end{equation}
o que comprova que o corte mínimo é o dual linear do fluxo máximo.

A representação esquemática dessa solução dual é ilustrada na Figura~\ref{fig:primal_dual_potenciais}.

\begin{figure}[!ht]
\centering
\begin{tikzpicture}[thick]
  \draw[cut S] (-0.8, -1.8) rectangle (2.8, 1.8);
  \node[cut label S] at (-0.6, 1.65) {$\mathbf{S} : p^*(u) = 1$};

  \draw[cut T] (5.4, -1.8) rectangle (9.0, 1.8);
  \node[cut label T] at (8.8, 1.65) {$\mathbf{T} : p^*(v) = 0$};

  \node[source node] (s) at (0.0, 0.0) {$s$};
  \node[s node internal] (a) at (2.0, 0.8) {$a$};
  \node[s node internal] (b) at (2.0, -0.8) {$b$};

  \node[t node internal] (c) at (6.2, 0.8) {$c$};
  \node[t node internal] (d) at (6.2, -0.8) {$d$};
  \node[sink node] (t) at (8.2, 0.0) {$t$};

  \draw[flow edge] (s) -- node[edge lbl above] {$d=0$} (a);
  \draw[flow edge] (s) -- node[edge lbl below] {$d=0$} (b);
  \draw[flow edge] (a) -- node[edge lbl left] {$d=0$} (b);

  \draw[flow edge] (c) -- node[edge lbl above] {$d=0$} (t);
  \draw[flow edge] (d) -- node[edge lbl below] {$d=0$} (t);
  \draw[flow edge] (c) -- node[edge lbl right] {$d=0$} (d);

  \draw[line width=1.5pt, dashed, cutviolet] (4.1, -1.85) -- (4.1, 1.85)
      node[above=2.5pt, solid, font=\scriptsize\bfseries, text=cutviolet!90!black, fill=white, inner sep=2.5pt, rounded corners=3pt, draw=cutviolet!70] {Corte Mínimo $(S,T)$};

  \draw[sat edge] (a) -- node[edge lbl above, text=cutviolet!90!black] {$\mathbf{d=1},\; f=c$} (c);
  \draw[sat edge] (b) -- node[edge lbl below, text=cutviolet!90!black] {$\mathbf{d=1},\; f=c$} (d);

  \node[font=\footnotesize, align=center, below=4pt] at (4.1, -2.0) {Custo dual: $\sum c(e)d(e) = c(a,c)\cdot 1 + c(b,d)\cdot 1 = \ccap(S,T)$. Pela dualidade forte: $\max |f| = \min \ccap(S,T)$.};
\end{tikzpicture}
\caption{Representação esquemática da solução dual no Teorema do Fluxo Máximo e Corte Mínimo: os potenciais nodais ótimos assumem valores binários $p^*(u) = 1$ no bloco S e $p^*(v) = 0$ no bloco $\mathbf{T}$. As variáveis duais de tensão $d(u, v) = \max\{0, p(u) - p(v)\}$ são ativadas ($d = 1$) exclusivamente sobre os arcos saturados que cruzam de S para T, anulando-se ($d = 0$) nas arestas internas ou reversas.}
\label{fig:primal_dual_potenciais}
\end{figure}


\section{O Teorema da Integralidade}

Para que problemas combinatórios discretos possam ser resolvidos via redes de fluxo, demonstra-se a integralidade da solução ótima.

\begin{teorema}[Integralidade de Ford-Fulkerson, 1956][teo:integralidade]
Se a função de capacidade $c(u,v)$ assume valores inteiros para todos os arcos $(u,v) \in A$, então existe um fluxo máximo $f^*$ onde $f^*(u,v) \in \mathbb{Z}$ para todo $(u,v) \in A$~\cite{ford1956maximal}.
\end{teorema}

\intuicao{dos Incrementos Discretos} Com fluxo inicial nulo $f_0 \equiv 0$ e capacidades inteiras, todas as capacidades residuais iniciais em $D_{f_0}$ são inteiras. A cada iteração de aumento, o gargalo residual $\delta$ ao longo de um caminho simples é um inteiro positivo. Como a adição e subtração de inteiros preserva a integralidade, o fluxo e as capacidades residuais mantêm-se inteiros em todas as etapas.

\begin{demonstracao}
Procedemos por indução no número $k$ de iterações de aumento de fluxo realizadas pelo método de Ford-Fulkerson:

\medskip
\noindent \textit{Base de indução ($k = 0$):} \\
O fluxo inicial é nulo: $f_0(u,v) = 0 \in \mathbb{Z}$ para todo $(u,v) \in A$. Logo, $f_0$ é um fluxo viável inteiro.

\medskip
\noindent \textit{Passo indutivo:} \\
Assuma que, após $k$ iterações, $f_k(u,v) \in \mathbb{Z}$ para todo $(u,v) \in A$. Como as capacidades $c(u,v)$ são inteiras, as capacidades residuais em $D_{f_k}$:
\begin{align*}
    c_{f_k}(u,v) &= c(u,v) - f_k(u,v) \in \mathbb{Z} \quad \text{(arcos diretos)} \\
    c_{f_k}(v,u) &= f_k(u,v) \in \mathbb{Z} \quad \text{(arcos reversos)}
\end{align*}
são inteiras e não negativas.

Se não existe caminho aumentante em $D_{f_k}$, o algoritmo termina e $f_k$ é o fluxo máximo, que é inteiro. Caso exista um caminho simples aumentante $P$ de $s$ a $t$ em $D_{f_k}$, o acréscimo de fluxo é:
$$\delta_k = \min_{(u,v) \in P} c_{f_k}(u,v)$$
Como $P$ é finito e $c_{f_k}(u,v) \in \mathbb{Z}^+$ para todo $(u,v) \in P$, tem-se $\delta_k \in \mathbb{Z}$ e $\delta_k \ge 1$.

O novo fluxo $f_{k+1}$ é obtido por:
$$f_{k+1}(u,v) = \begin{cases} f_k(u,v) + \delta_k, & \text{se } (u,v) \in P \text{ é direto} \\ f_k(u,v) - \delta_k, & \text{se } (v,u) \in P \text{ é reverso} \\ f_k(u,v), & \text{caso contrário} \end{cases}$$
Como $\mathbb{Z}$ é fechado sob adição e subtração, $f_{k+1}(u,v) \in \mathbb{Z}$ para todo $(u,v) \in A$.

Como as capacidades são finitas, a capacidade de qualquer corte $s-t$ é finita. Visto que cada iteração aumenta o fluxo em pelo menos uma unidade ($\delta_k \ge 1$), o algoritmo atinge a condição de término em número finito de iterações. Pelo Teorema~\ref{teo:mfmc}, o fluxo final $f^*$ é máximo e inteiro.
\end{demonstracao}

\implicacoes{e Reduções Combinatórias} Embora a formulação clássica de fluxo máximo seja um programa linear contínuo, a integralidade natural das soluções para capacidades inteiras elimina a necessidade de variáveis binárias explícitas. A solução ótima do problema contínuo coincide com um vértice inteiro do poliedro de fluxos. Essa propriedade viabiliza resolver problemas discretos (como emparelhamento bipartido e cortes de separação) em tempo estritamente polinomial.

\textbf{Observação construtiva:} Os algoritmos de aumento de fluxo (Ford-Fulkerson, Edmonds-Karp e Dinic) e de pré-fluxo (Push-Relabel) constroem soluções inteiras ao longo de todas as etapas quando operam sobre capacidades inteiras, partindo do fluxo nulo e aplicando apenas operações aritméticas em $\mathbb{Z}$. Do ponto de vista da programação linear, a integralidade dos vértices do poliedro decorre da \textbf{Total Unimodularidade} da matriz de incidência (demonstrada na Subseção~\ref{sec:tum_secao}).


